\section*{ЗАКЛЮЧЕНИЕ}
\addcontentsline{toc}{section}{ЗАКЛЮЧЕНИЕ}

В процессе выполнения выпускной квалификационной работы была разработана
компонентно-ориентированная программная система для построения интерактивных
приложений на языке Java. Разработанная система предоставляет программную
основу, которая позволяет создавать сцены, добавлять в них сущности и
компоненты, подключать системы обработки, выполнять пользовательские сценарии
поведения, обрабатывать ввод, отображать графические объекты и использовать
базовые физические взаимодействия.

При выполнении работы основной акцент был сделан на том, чтобы система не
сводилась к набору отдельных примеров, а имела общую архитектуру, пригодную для
расширения. Объекты сцены представлены как сущности, поведение и состояние
которых задаются через независимые компоненты. Обработка таких объектов
выполняется специализированными системами. Это позволило разделить разные
части логики приложения: управление сценой, рендеринг, обработку ввода,
выполнение пользовательских скриптов, отладочную визуализацию и физическую
обработку.

Основные результаты работы:
\begin{enumerate}
\item Проведен анализ предметной области. Рассмотрены интерактивные приложения
как класс программных систем, особенности их жизненного цикла, задачи
обработки пользовательского ввода, обновления состояния объектов и вывода
кадра. Изучены существующие программные средства построения интерактивных
приложений, включая низкоуровневые библиотеки, фреймворки и более крупные
готовые решения. На основе анализа обоснована целесообразность разработки
промежуточной программной основы, ориентированной на Java-разработчика.

\item Сформулированы требования к компонентно-ориентированной программной
системе. Определены назначение разработки, требования к архитектурной модели,
управлению приложением и сценами, графической подсистеме, пользовательскому
вводу, системе пользовательских скриптов, базовой физической подсистеме,
демонстрационным сценам и программному интерфейсу. Также описан порядок
использования системы разработчиком при создании собственной интерактивной
сцены.

\item Спроектирована архитектура программной системы. В качестве технологической
основы выбраны Java, LWJGL, OpenGL, GLSL, GLFW и STB. Определено внутреннее
устройство приложения, контекста выполнения, сцены, сущностей, компонентов,
систем обработки, запросов сущностей и ресурсов. Отдельно спроектированы
графическая подсистема, подсистема пользовательского ввода, подсистема
пользовательских скриптов, отладочная визуализация и физическая подсистема.

\item Реализованы основные классы и компоненты программной системы. В составе
реализации представлены ядро компонентной модели, классы управления
приложением и сценами, математические классы, компоненты трансформации,
камеры, графического представления, освещения, пользовательских скриптов,
физических тел и коллайдеров. Реализованы системы обработки, отвечающие за
выполнение скриптов, отрисовку сцены, отладочную визуализацию и базовую
физическую обработку.

\item Реализованы демонстрационные сцены, показывающие разные варианты
использования разработанной системы. Сцена \texttt{ShowcaseScene} демонстрирует
базовую трехмерную сцену, камеру, освещение и отладочные элементы.
\texttt{Interactive2DScene} показывает обработку пользовательского ввода и
работу интерактивных объектов. \texttt{TexturedLightingScene} демонстрирует
использование материалов, текстур и направленного освещения.
\texttt{PhysicsScene} проверяет совместную работу физической подсистемы,
рендеринга, скриптов, ввода и отладочной визуализации.

\item Проведено модульное, регрессионное и системное тестирование разработанной
программной системы. Были проверены ядро компонентной модели, подсистема
ввода, менеджер сцен и физическая подсистема. Системное тестирование
демонстрационных сцен подтвердило корректность создания окна приложения,
загрузки сцен, обработки пользовательских действий, отображения объектов,
работы освещения, текстур, физических тел и перезапуска сцены.
\end{enumerate}

Поставленная во введении цель достигнута. Разработанная программная система
позволяет создавать простые интерактивные приложения на Java без повторной
реализации базовых механизмов управления сценой, обработки ввода, организации
компонентов, рендеринга и физической обработки. Задачи, сформулированные в
начале работы, выполнены в полном объеме.

Практическая ценность результата заключается в том, что разработанная система
может использоваться как основа для учебных, демонстрационных и
исследовательских интерактивных приложений. Она также может быть полезна для
изучения внутреннего устройства компонентной архитектуры, так как основные
механизмы работы сцены, компонентов, систем обработки, графики и физики
реализованы явно и доступны для анализа и модификации.

Перспективами дальнейшего развития являются расширение графической
подсистемы, добавление новых типов материалов и источников света, поддержка
анимации, развитие физической модели, сохранение и загрузка описания сцен,
расширение набора готовых компонентов и создание более подробной
пользовательской документации. В дальнейшем на основе разработанного ядра
может быть создана более полнофункциональная программная платформа для
построения интерактивных приложений и игровых прототипов на языке Java.
